% Fichier généré automatiquement par tools/generate_corrections.py.
% Source : DYN/DYN-03-Inertie/41_Parallelepipede/41_Parallelepipede.tex
% Statut : complet (2/2 question(s) renseignée(s), 0 reprise(s)).
\begin{corrigebox}[Corrigé extrait de la banque]
\CorrigeQuestion{1}
On note $m_C$ la masse du cylindre (plein) et $m_P$ la masse du parallélépipède. On a alors $m=m_P-m_C$.
De plus, $\vect{OG_P}=\dfrac{a}{2}\vect{x}+\dfrac{b}{2}\vect{y}+\dfrac{c}{2}\vect{z}$
et $\vect{OG_C}=\dfrac{a}{3}\vect{x}+\dfrac{b}{2}\vect{y}+\dfrac{c}{2}\vect{z}$.

On a alors 
$m \vect{OG}=m_P \vect{OG_P}-m_C \vect{OG_C} = m_P \left(\dfrac{a}{2}\vect{x}+\dfrac{b}{2}\vect{y}+\dfrac{c}{2}\vect{z} \right)-m_C \left( \dfrac{a}{3}\vect{x}+\dfrac{b}{2}\vect{y}+\dfrac{c}{2}\vect{z}\right) $.

Par suite, 
$\vect{OG}=
\begin{pmatrix}
x_G \\ y_G \\ z_G
\end{pmatrix}_{\base{x}{y}{z}}
=
\begin{pmatrix}
\dfrac{a}{m_P-m_C} \left( \dfrac{m_P}{2} -  \dfrac{m_C}{3}\right) \\
b/2 \\
c/2
\end{pmatrix}_{\base{x}{y}{z}}
$.
\CorrigeQuestion{2}
Les plans $(G,\vect{x},\vect{y})$ et $(G,\vect{z},\vect{x})$ sont des plans de symétrie. 
On a donc $\inertie{G}{S}=\matinertie{A}{B}{C}{0}{0}{0}{\mathcal{B}}$.

\textbf{On déplace la matrice du parallélépipède rectangle en G.}

On a $\inertie{G_P}{P}=\matinertie{A_P}{B_P}{C_P}{0}{0}{0}{\mathcal{B}}$ et

$\vect{G_P G} = \vect{G_P O} + \vect{O G} = 
\begin{pmatrix}
-\dfrac{a}{2} \\
-\dfrac{b}{2} \\
-\dfrac{c}{2}
\end{pmatrix}_{\mathcal{B}} 
+
\begin{pmatrix}
\dfrac{a}{m_P-m_C} \left( \dfrac{m_P}{2} -  \dfrac{m_C}{3}\right) \\
b/2 \\
c/2
\end{pmatrix}_{\mathcal{B}}=
\begin{pmatrix}
\dfrac{a}{m_P-m_C} \left( \dfrac{m_P}{2} -  \dfrac{m_C}{3}\right) -\dfrac{a}{2}\\
0 \\
0
\end{pmatrix}_{\mathcal{B}}=
\begin{pmatrix}
\Delta_x\\
0 \\
0
\end{pmatrix}_{\mathcal{B}}$.

Ainsi, 
$\inertie{G}{P} = \inertie{G_P}{P} + m_P\begin{pmatrix}
0 & 0 & 0\\
0 & \Delta_x^2 & 0 \\
0 & 0 & \Delta_x^2
\end{pmatrix}_{\mathcal{B}}
= \matinertie{A_P}{B_P+m_P\Delta_x^2}{C_P+m_P\Delta_x^2}{0}{0}{0}{\mathcal{B}}
$

\textbf{On déplace la matrice du cylindre en G.}

De même $\inertie{G_C}{C}=\matinertie{A_C}{B_C}{A_C}{0}{0}{0}{\mathcal{B}}$ et

$\vect{G_C G} = \vect{G_C O} + \vect{O G} = 
\begin{pmatrix}
-\dfrac{a}{3} \\
-\dfrac{b}{2} \\
-\dfrac{c}{2}
\end{pmatrix}_{\mathcal{B}} 
+
\begin{pmatrix}
\dfrac{a}{m} \left( \dfrac{m_P}{2} -  \dfrac{m_C}{3}\right) \\
b/2 \\
c/2
\end{pmatrix}_{\mathcal{B}}=
\begin{pmatrix}
\dfrac{a}{m} \left( \dfrac{m_P}{2} -  \dfrac{m_C}{3}\right) -\dfrac{a}{3}\\
0 \\
0
\end{pmatrix}_{\mathcal{B}}=
\begin{pmatrix}
\Delta_x'\\
0 \\
0
\end{pmatrix}_{\mathcal{B}}$.

Ainsi, 
$\inertie{G}{C} = \inertie{G_C}{C} + m_C\begin{pmatrix}
0 & 0 & 0\\
0 & \Delta_x'^2 & 0 \\
0 & 0 & \Delta_x'^2
\end{pmatrix}_{\mathcal{B}}
= \matinertie{A_C}{B_C+m_C\Delta_x'^2}{A_C+m_C\Delta_x'^2}{0}{0}{0}{\mathcal{B}}
$.

\textbf{Bilan.}

Au final, $\inertie{G}{E} = \inertie{G}{P} - \inertie{G}{C}$ et 

$\inertie{G}{E} = \matinertie{A_P-A_C}{B_P+m_P\Delta_x^2-B_C-m_C\Delta_x'^2}{C_P+m_P\Delta_x^2-A_C-m_C\Delta_x'^2}{0}{0}{0}{\mathcal{B}} $
\end{corrigebox}
